% Retrieval Quality Under Permission-Shaped Filters: body of Section~\ref{sec:quality} (the \section
% heading and its label stay in main.tex). Labels defined here: sec:qual-audit, tab:qual-x0,
% sec:qual-qdrant, tab:qual-qdrant, sec:qual-pgvector, sec:qual-frontiers, sec:qual-size.
% Sources: evidence/kaggle-scale-run-2026-10-24.md (X0), evidence/x1-underretrieval-run1.md (X1 run
% 1), investigations/filtered-recall-gaps.md (mechanisms, with source links). X1 run 2, X2 and X3
% are pending and marked \tbd{}.

Security held in every configuration of Section~\ref{sec:e2e}; this section asks what the
permission filter cost. We first audit the enterprise scenario at common settings, then diagnose the
two largest losses with pre-registered experiments, and finally map recall against throughput and
filter size. Recall is permission-recall@$k$ against exact search over the readable set
(Section~\ref{sec:prob-metrics}).

\subsection{An audit at common settings}\label{sec:qual-audit}

\begin{table}[tbp]
\centering
\small
\begin{tabular}{@{}lrrrrrr@{}}
\toprule
 & \multicolumn{3}{c}{Four runs} & \multicolumn{2}{c}{Overlap@5, first run} & \\
\cmidrule(lr){2-4}\cmidrule(lr){5-6}
Store & Recall@5 & Recall@10 & Range@10 & judy & others & p50 (ms) \\
\midrule
FAISS (flat, exact) & 1.000 & 1.000 & 0.000 & 1.000 & 1.000 & 346--368 \\
Redis & 0.983 & 0.981 & 0.012 & 0.998 & 0.982 & 28--30 \\
Elasticsearch & 0.942 & 0.949 & 0.019 & 0.987 & 0.941 & 8--9 \\
OpenSearch & 0.942 & 0.935 & 0.005 & 0.998 & 0.941 & 7 \\
Weaviate & 0.899 & 0.902 & 0.021 & 1.000 & 0.876 & 7--8 \\
Chroma & 0.877 & 0.879 & 0.010 & 0.987 & 0.876 & 1432--1481 \\
LanceDB & 0.840 & 0.874 & 0.039 & 0.949 & 0.821 & 22--25 \\
pgvector, fast mode (2 runs) & 0.695 & 0.732 & 0.000 & 0.724 & 0.686 & 6.5--8.1 \\
pgvector, exact mode (2 runs) & 0.993 & 0.993 & 0.002 & & & 386--399 \\
Qdrant & 0.557 & 0.551 & 0.026 & 1.000 & 0.523 & 4.3--4.7 \\
\bottomrule
\end{tabular}
\caption{Retrieval quality in the enterprise scenario (Section~\ref{sec:e2e}), four runs of 1,105
questions (seeds 7, 11, 13, 17). Recall is the mean over runs; Range@10 is the largest minus the
smallest recall@10; p50 is the range of per-run medians. pgvector ran in one of two modes per run
(Section~\ref{sec:qual-pgvector}); each row averages the two runs in that mode. Overlap@5 is the
share of the exact top five that a store returned, by person type, from a re-analysis of the first
run (a fast-mode run for pgvector); ``judy'' are the least-privileged people (624 readable chunks).
pgvector with OPA follows pgvector's mode in each run. Settings: Appendix~\ref{app:settings}.}
\label{tab:qual-x0}
\end{table}

Table~\ref{tab:qual-x0} shows the spread. At settings a team would plausibly ship, mean recall@10
ranges from 0.55 for Qdrant to 0.98 for Redis, and every store's loss is concentrated where the
readable set is far from the question. Most stores repeat closely: across four runs with fresh
documents and questions, no store other than pgvector moved its recall@10 by more than 0.04. Two
failures stand out. pgvector returned \emph{nothing} to some of the least-privileged people,
although each could read 624 chunks: ten in the first run, eight on red-team prompts and two on
questions about documents they could not read, and seven in another; the layer, before the change
of Section~\ref{sec:sys-complete}, called these empty answers complete. In the other two runs the
same configuration searched exactly, at fifty times the latency, and left one person without results
in one of them. Qdrant never returned a
short answer, yet for ordinary employees it returned the right documents only about half the time
(overlap@5 0.52 in the first run), while for the least-privileged people it was perfect. On targeted
questions in that run, 59.6\% of Qdrant's top five were shared FiQA passages where exact search
returned none, and the model found the target document for only 39.9\% of targeted questions
(33.8\% to 39.9\% over the four runs).

The two patterns are opposite in kind. pgvector \emph{under-retrieves}: it returns fewer than $k$
hits, which a count can detect. Qdrant \emph{mis-retrieves}: it returns $k$ hits, the wrong ones,
which only an exact baseline reveals. We pre-registered hypotheses for both before running a
diagnosis grid on Kaggle (X1; the pre-registration and its amendments are released).

\subsection{Qdrant: an islanded permitted subgraph}\label{sec:qual-qdrant}

Qdrant plans filtered search per segment from estimated filter cardinality: a filter estimated below
a threshold (about 3,333 points for 768-dimensional vectors at the default) runs as a plain, exact
scan over the payload index, and a broader one runs as filtered HNSW traversal, which scores and
expands only neighbors that pass the filter~\cite{qdrant2026search}. Extra graph edges that keep a
filtered subgraph connected are built per payload value, and only for values held by fewer points
than a cap of a few percent of the segment~\cite{couzon2026filtered}. In the scenario every tenant
value (56\% of points) and
the employee and contractor groups (65\% and 61\%) exceed the cap. Near a question about company
documents, only one of the twenty near-identical tenant copies passes the filter, so the permitted
company region is a sparse island in the graph, while the shared passages, all permitted, form a
large connected region that the traversal enters and does not leave. The least-privileged people
escape the trap because their filter is small enough to be planned as an exact scan.

\begin{table}[tbp]
\centering
\small
\begin{tabular}{@{}llrrr@{}}
\toprule
Layout & Search setting & Recall@10 & Range over builds & p50 (ms) \\
\midrule
As run (one collection) & \texttt{hnsw\_ef} 128 (as run) & 0.619 & 0.062 & 4.5 \\
 & \texttt{hnsw\_ef} 64 & 0.571 & 0.049 & 3.9 \\
 & \texttt{hnsw\_ef} 512 & 0.709 & 0.101 & 7.7 \\
 & ACORN, max selectivity 0.4 or 1.0 & 0.761 & 0.151 & 4.8 \\
 & exact & 1.000 & n/a & 35.4 \\
Tenant field marked \texttt{is\_tenant} & as run & 0.656 & 0.031 & 4.3 \\
 & ACORN & 0.830 & 0.038 & 4.5 \\
Per-tenant graphs (\texttt{is\_tenant}, $m = 0$) & as run & 0.943 & 0.009 & 3.8 \\
Shared passages in a separate collection & as run & 0.948 & 0.004 & 8.2 \\
\bottomrule
\end{tabular}
\caption{Qdrant 1.19.2 under the enterprise scenario's permission filters (X1, 1,489 questions,
three index builds per layout, means over builds). ACORN at the default gate (0.4) and at 1.0 gave
identical results. \tbd{one-segment layout after the run-2 fix.}}
\label{tab:qual-qdrant}
\end{table}

The pre-registered predictions for this hypothesis all held (Table~\ref{tab:qual-qdrant}). Recall@10
for the least-privileged people was 1.000, against 0.273 (95\% CI 0.242 to 0.305) for company
documents among everyone else; 98.6\% of the hits that did not count were shared passages; and two
interventions that reconnect the permitted region restored recall. ACORN, which explores neighbors
of filtered-out neighbors~\cite{patel2024acorn}, lifted recall@10 from 0.619 to 0.761 at an 8\%
latency cost, and per-tenant graphs~\cite{qdrant2026multitenancy} lifted it to 0.943 while also
being the fastest configuration.
Putting the shared passages in their own collection, queried separately, reached 0.948. Raising
\texttt{hnsw\_ef} from 128 to 512 added 0.090 at 1.7 times the latency, within the band we had
pre-registered as inconclusive, and did not close the gap, as expected for a traversal trapped in
the wrong region.

One prediction failed. We expected ACORN's default gate, which engages when the estimated
selectivity is at most 0.4, to engage for some people and not others, because the planner's
independence estimate for a typical employee sits near 0.4. It engaged for everyone, and the gate
at 1.0 gave identical recall. The paper therefore claims that ACORN helps but does not suffice under
this layout, not that its gate misfires. A third hypothesis, that recall varies with the segment
layout, remains open: recall@10 varied by 0.062 between three builds of the same data, but our
attempt to force a single segment did not take effect. \tbd{run 2}

For practitioners the result is a layout rule, not a parameter: when permission groups or tenants
cover large, semantically interleaved parts of a collection, give each tenant its own graph, or keep
widely shared content apart, before reaching for search-time effort.

\subsection{pgvector: the iterative scan's budget}\label{sec:qual-pgvector}

pgvector 0.8 answers a filtered HNSW query by scanning the graph in order of distance and applying
the filter to each row the scan emits; with iterative scans enabled, it continues past
\texttt{hnsw.ef\_search} until enough rows pass the filter or a limit is
reached~\cite{pgvector2026pgvector}. The limit is \texttt{hnsw.max\_scan\_tuples} (20,000 by default)
or a memory bound of \texttt{hnsw.scan\_mem\_multiplier} times \texttt{work\_mem}. A person whose
readable set is small and far from the question can exhaust the budget before any readable row is
emitted, and the query then returns nothing, which is indistinguishable from having nothing to read.

The pre-registered test of this mechanism held. Every empty answer stopped after more than half of
the budget had been spent, at a median of 20,083 rows emitted, and every setting that raised the
tuple limit to 100,000 or more removed the empty answers entirely, while raising only the memory
bound (eight times the multiplier, or 64\,MB of \texttt{work\_mem}) at the default tuple limit did
not. \tbd{X1 run 2: absolute counts; the magnitude of the X0
loss (recall@5 0.691 in the scenario run against 0.998 in the first diagnosis run) and its
explanation from investigations/x0-x1-pgvector-gap.md; ef\_search and build effects; exact
baseline.}

\subsection{Recall against throughput}\label{sec:qual-frontiers}

The enterprise scenario has one permission layout. To separate the filter's selectivity from its
correlation with the question, X2 adds a corpus whose layouts we control: the first million
passages of MS MARCO~\cite{bajaj2016ms} from BEIR, embedded with bge-small-en-v1.5 (384 dimensions,
so that ten stores fit beside the harness on a 16\,GB runner). Every passage belongs to one team of
a hierarchy of 100 divisions, 10 departments each and 10 teams each, assigned either at random
(\emph{uncorrelated}) or by hierarchical $k$-means over the embeddings (\emph{clustered}, so groups
are topical), five seeds of each. A subject receives divisions, departments and teams until it may
read a target share of the corpus, from 0.06\% to 100\%, and its filter is either one disjunction over
its groups or a tenant clause conjoined with a disjunction over a grant twice as large, a shape in
which each clause alone passes far more than both. For clustered layouts the questions are drawn
either from those whose 100 nearest passages are at least half readable (\emph{positively
correlated}) or from those with nothing readable among their 1,000 nearest (\emph{negatively
correlated}). For every question and subject we compute, from exact scores, the global selectivity,
the local selectivity among the 100 and 1,000 nearest passages, the rank of the first readable
passage, ACORN's query correlation~\cite{patel2024acorn}, the global-local selectivity of Amanbayev
et al.~\cite{amanbayev2026filtered} and the over-fetch estimate of HCBGen~\cite{lim2026revisiting}.

Ten approximate stores run against FAISS over a flat index: pgvector, Qdrant, Elasticsearch,
OpenSearch, Weaviate, Redis, Milvus, MongoDB, LanceDB and Chroma, each at its default, at five levels
of its search-effort parameter, with its store-specific options (pgvector's raised scan limits and
strict order, Qdrant's ACORN and per-tenant graphs, Weaviate's filter strategies, Redis's hybrid
policies, Elasticsearch's quantized index types) and with its own exact search where it has one, in
three fresh builds per store. Each pgvector search is planned with its own values, so the plan-cache
effect of Section~\ref{sec:qual-pgvector} cannot arise, and the share of searches whose plan uses the
HNSW index is recorded. Throughput is measured end to end through the layer with one client thread
and compared only within a job; recall is compared across jobs. The same harness runs the enterprise
scenario, whose as-run settings must reproduce the audit's recall within 0.03 for seven of nine
stores. We pre-registered four hypotheses: that negatively correlated questions lose at least 0.05
of recall@10 against uncorrelated ones at defaults in at least six stores (H2.1); that local
selectivity predicts recall better than global selectivity (H2.2); that raising search effort
narrows but does not close that gap (H2.3); and that intermediate selectivities are harder than very
small ones, which stores answer by brute force (H2.4).
\tbd{X2 full run (GitHub Actions run 40921769714): frontiers per store and regime (figure);
verdicts on H2.1 to H2.4 with intervals; recall by measured axis; validity checks V1 to V5.}

\subsection{Filter size}\label{sec:qual-size}

A person in many directory groups turns the permission filter into a long disjunction: the filter of
Section~\ref{sec:sys-contracts} holds one value per group, and enterprise users can belong to
thousands. We vary only that length. Each of five seeds builds a corpus of 100,000 vectors of 384
dimensions, the 57,638 FiQA passages embedded as in Section~\ref{sec:side-setup} and 42,362 synthetic
vectors around 200 seeded cluster centers, and makes a fixed 10\% of it readable at every level: a
readable document holds one group of each level, and the subject at level $G$ holds exactly that
level's $G$ groups, for $G$ from 1 to 40,000. The readable set, and so the exact answer, is the same
at every level; only the filter grows. Twelve stores run at default index settings, with FAISS
over a flat index and the in-memory backend as controls, and each level is timed over 30 questions
with $k = 10$ in three passes after a warm-up, a retrieval being allowed 120\,s. We record whether
each store's translator returns and how fast, whether the store accepts the filter or refuses it,
and with which error, the latency percentiles, and recall@10 against exact search over the readable
set; for pgvector also the plan. The limits the stores document that could bind are a 65,536-term
cap per \texttt{terms} query in Elasticsearch and OpenSearch, PostgreSQL's 65,535 bind parameters,
SQLite's 32,766 host parameters and expression depth of 1,000 under Chroma, and the request sizes of
MongoDB (16\,MB) and Qdrant (32\,MB); the layer's Cassandra adapter falls back to one unfiltered
query filtered locally beyond 254 groups, as documented.

We pre-registered that every translator returns within a second at every level; that at 40,000
groups Elasticsearch, OpenSearch, Qdrant, pgvector, LanceDB, Cassandra, FAISS and the in-memory
backend accept the filter and Chroma refuses it; that pgvector's median latency grows at least
fiftyfold from one group to 40,000, because above 100 values the layer sends an array test that the
GIN index cannot serve, while the controls' grows at most twofold; and that recall does not depend
on the filter's length in any approximate store that accepts it (paired difference within
$\pm 0.02$).
\tbd{X3 full run (GitHub Actions run 40940385704): acceptance per store and level with error
classes; latency against $G$; recall differences with intervals; pgvector's plans; Chroma's
process ended with a segmentation fault (signal 11) in every seed it reached, rather than refusing
the filter with an error, to be confirmed from the logs.}
